\chapter{Introdução}
\label{cap:introducao}

\section{Contexto}

Na relatividade geral (RG), uma onda gravitacional no vácuo propaga-se
com a velocidade da luz e possui dois estados de polarização, os modos
``$+$'' e ``$\times$''. Ambas as propriedades decorrem de o gráviton,
na linguagem de teoria de campos, ser uma partícula de spin dois sem
massa. Uma massa não nula muda as duas coisas ao mesmo tempo: a relação
de dispersão passa a depender da frequência e surgem estados de
polarização adicionais. Por isso, a propagação e a polarização das ondas
gravitacionais são, desde os anos 1970, instrumentos naturais para testar
a gravitação \autocite{eardley1973,will2014}.

Esses testes deixaram de ser apenas propostas. Desde a primeira detecção
direta de uma coalescência de buracos negros \autocite{gw150914}, as
colaborações LIGO, Virgo e KAGRA (LVK) restringem a massa do gráviton
pela dispersão acumulada ao longo da propagação; o limite mais recente
é $m_g\leq1{,}92\times10^{-23}\,\mathrm{eV}/c^2$ a 90\% de credibilidade
\autocite{lvk2026}. Em 2023, redes de temporização de pulsares (PTAs)
reportaram evidência de um fundo estocástico de ondas gravitacionais na
faixa de nanohertz, identificado pela correlação angular de
Hellings--Downs entre os resíduos de diferentes pulsares
\autocite{agazie2023,hellings1983}. Esses dados já foram usados para
limitar a massa do gráviton \autocite{wu2024} e para buscar polarizações
transversais adicionais \autocite{nanograv2024}.

\section{O trabalho de partida}

Esta dissertação parte do Trabalho de Graduação de Wayne Leonardo Silva
de Paula, defendido no ITA em 2003 \autocite{tg2003}, cujo núcleo foi
publicado no ano seguinte por de Paula, Miranda e Marinho
\autocite{depaula2004}. O TG estudou a teoria bimétrica de Visser
\autocite{visser1998}, na qual a massa do gráviton é introduzida por um
termo que depende de uma métrica de fundo não dinâmica. Linearizando as
equações de campo em torno de Minkowski e aplicando o formalismo de
Newman--Penrose (NP) e a classificação E(2) de Eardley, Lee e Lightman
\autocite{eardley1973}, o trabalho obteve três resultados:
\begin{enumerate}
 \item as ondas planas fracas da teoria pertencem à classe~II$_6$, isto é,
 admitem até seis estados de polarização, com conteúdo de helicidades
 $0$, $\pm1$ e $\pm2$;
 \item expressões explícitas para as amplitudes NP
 $\Psi_2,\Psi_3,\Psi_4$ e $\Phi_{22}$ em termos das componentes da
 perturbação métrica, e a matriz de marés $E_{ij}$ que atuaria sobre
 um detector;
 \item uma comparação entre bandas de frequência: para amplitudes métricas
 comparáveis, os modos adicionais seriam desprezíveis na banda dos
 interferômetros terrestres e tornar-se-iam comparáveis aos modos da RG
 em frequências próximas de $10^{-7}\,\mathrm{Hz}$. O TG concluiu que
 detectores de baixa frequência seriam os mais adequados para
 discriminar a teoria massiva.
\end{enumerate}

\section{O que mudou em vinte anos}

Três desenvolvimentos posteriores tornam oportuno revisitar esses
resultados.

\paragraph{Teoria.} Sabe-se hoje que, entre os termos de massa
quadráticos em torno de Minkowski, apenas a combinação de Fierz--Pauli
\autocite{fierz1939} propaga exatamente cinco graus de liberdade sem
fantasmas. Qualquer outra combinação propaga um sexto modo escalar com
energia cinética de sinal errado \autocite{hinterbichler2012,park2010}.
O termo de Visser não é de Fierz--Pauli; o próprio autor observa que
essa escolha é deliberada, para evitar a descontinuidade de van
Dam--Veltman--Zakharov (vDVZ) no limite de massa nula
\autocite{visser1998,vandam1970,zakharov1970}. A contagem de seis
polarizações do TG deve, portanto, ser reinterpretada à luz desse
conhecimento. Também se compreende melhor a estrutura não linear:
o fantasma de Boulware--Deser, sua eliminação na teoria de de Rham,
Gabadadze e Tolley (dRGT) e o mecanismo de Vainshtein
\autocite{boulware1972,derham2011,vainshtein1972,babichev2013}.

\paragraph{Limites.} O limite do Sistema Solar usado em 2004 era
$m_g<4{,}4\times10^{-22}\,\mathrm{eV}/c^2$. Os limites atuais são cerca
de vinte a cinquenta vezes menores. Como os efeitos cinemáticos da massa
são controlados pela razão $(f_g/f)^2$, em que $f_g=m_gc^2/h$ é a
frequência de corte, a banda em que os modos adicionais deixam de ser
suprimidos desloca-se de $10^{-7}\,\mathrm{Hz}$ para poucos nanohertz.
Essa é precisamente a banda das PTAs.

\paragraph{Formalismo.} A classificação E(2) foi construída para ondas
nulas ou quase nulas. Ela pressupõe que a diferença entre a velocidade
da onda e a da luz seja pequena. No regime em que os modos adicionais
seriam relevantes, $f\sim f_g$, essa hipótese deixa de valer. O próprio
exemplo numérico de baixa frequência do artigo de 2004 corresponde a uma
velocidade de grupo de cerca de um quarto da velocidade da luz. Uma
revisão do resultado exige, portanto, um tratamento exato em $f/f_g$
\autocite{hyun2019}.

\section{Questões e objetivos}

A dissertação busca responder três questões encadeadas.

\begin{enumerate}
 \item[Q1.] \emph{O que permanece correto nos resultados do TG quando a
 análise é refeita sem a aproximação quase nula, e como eles se comparam
 à teoria linear consistente de Fierz--Pauli?}
 \item[Q2.] \emph{Se os modos adicionais só podem ter efeito apreciável
 na banda de nanohertz, como eles aparecem nas correlações entre
 pulsares medidas por uma PTA?}
 \item[Q3.] \emph{Quanto uma PTA pequena consegue efetivamente aprender
 sobre a massa e sobre um modo escalar, e que cuidados de análise esses
 testes exigem?}
\end{enumerate}

Os objetivos específicos correspondentes são:
\begin{enumerate}
 \item reconstruir a matriz de marés de uma onda plana massiva sem supor
 propagação nula; contar as amplitudes independentes no modelo de Visser
 e em Fierz--Pauli; identificar a origem física do sexto modo;
 \item delimitar o domínio de validade das expressões NP do TG e da
 classificação E(2) para ondas massivas, e rever a comparação entre
 bandas de frequência com os limites atuais;
 \item derivar a resposta de temporização de pulsares a uma onda massiva,
 incluindo o termo do pulsar e o movimento dos observadores, e calcular
 as correlações angulares dos setores tensorial e de helicidade zero;
 \item usar simulações controladas, já executadas, para estimar a
 informação disponível sobre a massa e sobre o modo escalar em redes
 pequenas, e identificar aproximações de análise que distorcem as
 conclusões.
\end{enumerate}

\section{Contribuições e delimitação}

As contribuições desta dissertação são as seguintes.
\begin{itemize}
 \item Uma revisão exata, em $f/f_g$, da cinemática de polarização do TG.
 Confirma-se que o modelo linear de Visser possui seis amplitudes de maré
 independentes. Mostra-se que a sexta é o modo escalar associado ao traço
 da perturbação, que na teoria linear não Fierz--Pauli é um fantasma. Em
 Fierz--Pauli, restam cinco amplitudes, e as deformações escalares
 transversal e longitudinal ficam vinculadas por
 $p_l/p_b=-(f_g/f)^2$.
 \item A identificação dos limites de validade das amplitudes NP do TG:
 o erro da leitura quase nula de $\Psi_4$ é
 $1-(1+\beta)^2/4$, e as classes E(2) não são invariantes para ondas
 massivas. Registram-se também uma ambiguidade de fator dois na
 normalização da massa e uma convenção de sinal da curvatura.
 \item A demonstração de que a conclusão qualitativa do TG sobre as bandas
 de frequência se mantém, mas depende da hipótese de amplitudes métricas
 comparáveis. Com os limites atuais, a razão entre marés não tensoriais
 e tensoriais só é de ordem um na banda das PTAs.
 \item A resposta de PTA para os setores tensorial e de helicidade zero
 de Fierz--Pauli, com termo do pulsar e movimento dos observadores. Ela
 é expressa diretamente pela matriz de marés, o que liga o objeto
 calculado no TG ao observável de temporização. Mostra-se que, no limiar
 $f=f_g$, todas as helicidades produzem a mesma correlação angular por
 estado, e explica-se essa degenerescência pela simetria de rotação do
 referencial de repouso da onda. Identifica-se também uma discrepância
 numa expressão publicada da correlação escalar.
 \item Uma avaliação honesta, a partir de campanhas sintéticas já
 executadas, do pouco que redes pequenas aprendem sobre a massa. Mostra-se
 que a aproximação normal usual para estimadores de correlação multiplica
 por cerca de dez a taxa de falsas detecções de polarização escalar.
\end{itemize}

A dissertação não deduz a emissão de ondas por fontes astrofísicas em uma
teoria massiva completa, não resolve o mecanismo de Vainshtein e não
analisa dados observacionais de PTAs. As amplitudes dos modos adicionais
são tratadas como parâmetros fenomenológicos. As simulações usam redes
pequenas, ruído conhecido e sinal fraco; seus resultados não são
restrições observacionais. Ao longo do texto, essas limitações são
discutidas no ponto em que afetam cada conclusão.

\section{Organização}

O Capítulo~\ref{cap:fundamentos} revisa as polarizações na RG e em
teorias métricas gerais, a gravitação massiva, os limites observacionais
e a detecção por PTAs. O Capítulo~\ref{cap:polarizacoes} contém a revisão
do TG: matriz de marés exata, contagem de modos em Visser e em
Fierz--Pauli, validade do formalismo NP e comparação entre bandas. O
Capítulo~\ref{cap:pta} leva esses resultados às correlações entre
pulsares. O Capítulo~\ref{cap:simulacoes} apresenta os experimentos
sintéticos e o que eles ensinam sobre a capacidade de uma PTA pequena.
O Capítulo~\ref{cap:conclusoes} responde às questões Q1--Q3 e indica os
próximos passos. O Apêndice~\ref{ap:convencoes} reúne convenções, e o
Apêndice~\ref{ap:reproducao} relaciona cada resultado aos programas que o
produzem.
